% 【专题研讨】所属：第4章（由 chapters/revised/04-sources.tex \input 引入）
\minitopic{专题研讨：记忆保存、来源监控与归纳规则}记忆与归纳都把主体带到“当前未直接给出”的内容。记忆让过去的信息在现在继续发挥理由作用，归纳让已观察样本支持尚未观察情形。二者都不是机械复制：记忆会重构，归纳依赖分类和背景假设。因而核心问题不是“是否百分之百可靠”，而是何种保存、生成和校验机制能使跨时间推理获得可撤销的认识资格。

\minitopic{记忆只保存证成吗}保存论认为，记忆通常把过去已经获得的内容与证成延续到现在，而不凭空创造新的认识资格。你昨天亲眼确认门已锁，今天记得此事，即使不再重现视觉细节，仍可依赖记忆。若昨天只是随口猜测，今天鲜明想起“门锁了”，记忆不会把原猜测升级为知识。保存论很好地解释了记忆对来源历史的依赖。 但记忆似乎也能产生新确证。你分别记得甲比乙高、乙比丙高，今天第一次推出甲比丙高；或多个模糊片段在新线索帮助下形成对事件结构的理解。这里新结论并非过去已经相信。生成论可说，记忆储存的材料在当前推理中共同提供新理由。较稳妥的区分是：记忆内容的资格通常受原始取得约束，记忆材料的重新组合则可能生成对新命题的证成。 保存还涉及“被遗忘的证据”问题。学生通过可靠教材学到某历史日期，数年后只记得日期，不记得出处。外在主义可以强调记忆链仍正常保存知识；严格内在主义担心主体无法回答任何来源质疑。一个中间方案要求主体保留最低来源标记，例如“这是受过检验的学习内容，而非梦境或猜测”，并在具体反证出现时愿意查询。日常知识无需随身携带完整档案，却不能与来源完全脱钩。

\minitopic{来源监控与记忆污染}记忆错误常不是内容完全消失，而是来源错配。你记得一句观点，却把自己想到、同事说过和文章读到混为一谈；目击者在新闻报道后把报道细节纳入现场记忆；研究者多次打开修改后的数据表，误以为其中一列从实验当天就存在。来源监控评价主体能否区分信息从何而来、在何时获得、经过哪些转述。 鲜明度不是可靠来源标志。反复讲述、图像想象和他人确认都能增加主观流畅，却也可能让错误内容更熟悉。群体一致同样可能来自共同污染：证人彼此讨论后形成趋同叙述，不等于获得多份独立证据。因此调查程序重视最早陈述、独立访谈和提问措辞。程序设计不是对“真实记忆”的外在干预，而是保存证据独立性的组成部分。 数字记录既能校验也能改写记忆。照片时间戳、版本历史和录音可提供外部约束，但剪辑、缺帧、自动摘要和错误元数据也会制造新的来源问题。记录不是天然优于人脑；应比较两条证据链的完整性、可篡改性和独立支持。最佳判断往往来自异质来源交叉：原始日志、设备记录、当事人早期笔记与互不接触的证言。

\minitopic{自传记忆与自我知识}自传记忆不只储存事实，也维持“这是发生在我身上的事”的连续性。主体对自己的经历通常有第一人称权威，但这种权威不是历史档案式不可错。情绪、家庭叙事和后来的自我理解会改变事件的突出部分。一个人可以无可置疑地知道自己现在有某种痛苦回忆，却错误判断童年事件的具体原因、日期或他人意图。 因此应区分现象报告、事件主张和解释主张。“我现在仿佛记得一间红色房间”首先报告当前记忆现象；“房间当时确实是红色”是关于过去的事实主张；“它说明家人故意欺骗我”又是因果与意图解释。三者需要不同证据。尊重第一人称经验不要求把所有历史和解释层面的结论都免于核验。 记忆争议也有伦理边界。追求事实不授权强迫当事人反复重演创伤，机构可以在保护隐私与保留证据之间采用分级访问。认识上的谨慎语言尤其重要：“记录未能确认”不同于“事情没有发生”，“记忆可能受影响”不同于“当事人在说谎”。前者保留证据状态，后者对人格与动机作出更强判断。

\minitopic{归纳问题不等于不许预测}休谟指出，从已观察情形到未观察情形的推理不能由演绎保证；若以“自然过去一贯”证明自然未来仍一贯，又使用了同类归纳。\parencite{hume2000} 这个挑战不是说每次预测都不合理，而是要求说明归纳资格来自何处。把问题回答为“科学一直很成功”只增加了新的历史样本，未提供非循环的最终证明。 可能回应包括：把归纳视为无法再证明但可被实践辩护的基本规则；用概率一致性和准确性规范约束置信度；强调特定领域中因果机制、稳定条件和模型检验；或接受怀疑结论，放弃寻求演绎式基础。各方案解释不同层面。概率规则能说明信念怎样保持一致，却不独自指定先验和模型；因果知识能支持某项外推，却仍需说明机制在目标环境中持续。 科学归纳很少只靠“过去像未来”。疫苗试验向更大人群外推，要判断样本选择、剂量机制、年龄差异和流行变异；天文预测依赖动力学模型，而非简单频率；机器学习从训练集到部署环境，则要求输入分布与目标关系保持足够稳定。归纳证成是样本、分类、机制与环境不变性共同作用的结果。

\minitopic{古德曼问题与分类选择}即使承认过去支持未来，仍有“按何种谓词投射”的问题。观察到的翡翠都绿色，既符合“绿色”，也可符合一个人为规定的谓词：在某日期前为绿色、此后为蓝色。两种概括都与旧样本一致，却给未来相反预测。问题表明，数据本身不唯一决定分类规则；语言习惯、理论嵌入、机制与项目性共同限制可投射谓词。 不能只说怪异谓词“看起来不自然”，因为自然性正是要解释的。较好的分析问：分类是否在不同测量下稳定，是否参与成功因果解释，能否支持多种独立预测，是否为了适配既有数据而事后构造。一个模型用复杂条件精确拟合全部旧样本，却对轻微时间改变给出跳跃预测，通常缺少稳健性。奥卡姆式简洁是线索，但简洁也依赖表示语言。 这对算法泛化有直接意义。训练准确率只说明模型适配已见数据；若模型抓住背景水印而非病理特征，部署到新医院就失败。反事实测试、跨环境验证和特征消融可帮助识别模型实际投射的分类。哲学上的“新谜题”因此转化为研究设计问题：哪些不变关系值得外推？

\minitopic{先验与经验的协作}先验确证不以某次特定感官经验为最终根据，不等于主体从未通过经验学会概念。理解“所有正方形有四条边”需要语言和图形训练，但一旦掌握概念，对命题的理由似乎来自概念关系。经验可以是获得理解的条件，却不是每次确证的证据。这一区分防止把先验误解为出生时自带的事实清单。 所谓直觉也不能自动取得先验权威。哲学案例判断可能受措辞、文化背景和理论训练影响；数学直观也可被形式证明纠正。先验探究同样需要反例、概念澄清和推理审计。可修正性不等于经验性：一个主体可以发现自己误解了概念或漏掉推理步骤，而不是通过新观察直接反驳逻辑关系。 先验原则与经验材料在科学中持续协作。概率公理约束置信度结构，实际概率值来自数据与模型；逻辑规则约束推导，是否应用某一模型取决于观测；价值与决策原则规定怎样权衡，损益估计来自世界。把二者分成互不接触的领域，会看不见真实推理的层级。

\minitopic{比较视野：温故、积习与类推}《论语》中“温故”可被理解为通过重访旧材料获得新的把握，而不只是机械复述。\parencite{analects2003} 这为保存与生成的区分提供比较问题：复习何时只是维持，何时因重新组织而产生理解？但它首先属于学习与人格养成的实践语境，不能直接当作当代记忆模块理论。 先秦论辩中的类推与辨类也提示，外推依赖“哪些相似性相关”。孟子和荀子的不同论证都频繁从具体情形引向较一般判断，却对人性、环境与养成的机制有不同设定。\parencite{mencius2008,xunzi2014} 比较研究不应只数相似例句，而应重建类比的桥梁：被比较对象共享什么结构，差异为何不破坏推理，结论适用到何处。

\minitopic{综合案例：城市洪水预测}某城市用过去二十年的降雨、地形与排水数据训练洪水模型。模型在历史回测中表现优秀，却在一次新型极端降雨中低估积水。工程团队发现，过去五年城市新增大量不透水地面，若干传感器也曾在高水位饱和。请分别分析记忆与归纳问题。 数据档案是城市的制度记忆，但版本、缺测和传感器上限影响来源监控；回测复现同一历史关系，不能证明部署环境未改变。新增地面改变因果机制，极端降雨超出样本范围，传感器饱和又使“真实峰值”没有被正确记录。失败可能来自错误保存、分类不当、分布变化或机制改变，修复方案随诊断不同。 团队不应简单删除旧模型。应保存故障时间线，重建原始传感器范围，引入水文机制，分区域测试不透水率影响，并报告在未见强度下的不确定性。模型仍可为一般天气提供知识，同时不再支持对极端事件的强断言。这个案例说明，合理归纳的标志不是从不失败，而是能定位外推边界并从失败中修正规则。

\begin{tcolorbox}[breakable,colback=white,colframe=blue!30!black,title={专题研读任务},fonttitle=\bfseries]
选择一个依赖历史记录的现实预测，分别建立“记忆链”和“归纳链”：前者追踪数据如何保存、改写和认证，后者追踪样本如何被分类并外推。找出一个只破坏保存的风险和一个只破坏外推的风险。
\end{tcolorbox}

\index{记忆保存论}
\index{来源监控}
\index{自传记忆}
\index{归纳问题}
\index{可投射性}
\index{先验知识}
